\section{Ordinary instanton theory and the handle units}
\label{sec:ordinary}

We write $Y_j=S^3_j(K)$, including $Y_\infty=S^3$, and write
$A\subset Y_0$ for the capped Seifert surface.  On $Y_0$ the determinant
line evaluates oddly on $A$; on $Y_{\pm1},Y_{\pm2}$ and $Y_\infty$ it is
trivializable.  All gauge transformations have determinant one.  In
particular, fixing a determinant means fixing its connection and its
transition identifications, not passing to the full group of projective
gauge transformations.

We establish the following ordinary input.  The word ``unit'' means an
isomorphism on instanton homology over $\F_2$; it does not assert that a
particular matrix on the chain group is invertible.

\begin{theorem}[Ordinary handle units]\label{thm:ordinary:units}
Let $C(Y_j;\Z)$ be the irreducible, unframed instanton complex with the
determinants just specified.  Compatible small perturbations and coherent
orientations can be chosen so that the rigid counts define integral
chain maps, and the rigid and interval counts have the following
properties over $\F_2$.
\begin{enumerate}[label=\textup{(\roman*)}]
\item Each rigid meridional handle map $Y_r\longrightarrow Y_{r+1}$,
for $r=0,1$, is a unit whenever its determinant has the specified end
restrictions.  The rigid count on their composite is their composition
up to chain homotopy.  Denote this composite by
$f_+:C(Y_0)\longrightarrow C(Y_2)$.
\item On the path $2\longrightarrow\infty\longrightarrow0$, let
$F_l,A$ be the capped Seifert classes of squares $-2,0$, respectively,
with $F_l\cdot A=0$, and let $S_l=A-F_l$.  Choose the base determinant
$c(F_l,A)=(0,1)$ and the two determinants
$c+e\PD(S_l)$, $e=0,1$, identified off $S_l$.  The sum modulo two of
the interval counts of actual index $-1$ is a unit
$g_+:C(Y_2;\F_2)\longrightarrow C(Y_0;\F_2)$.
\item The mirror construction, read backwards, gives units
$f_-:C(Y_{-2})\longrightarrow C(Y_0)$ on the rigid path through $-1$,
and $g_-:C(Y_0;\F_2)\longrightarrow C(Y_{-2};\F_2)$ on
$0\longrightarrow\infty\longrightarrow-2$.  For $g_-$ the base
evaluations are $(1,0)$ and the second determinant differs by the
Poincar\'e dual of the intervening $(-2)$-sphere.
\end{enumerate}
For the summed interval maps, each summand has an integral orientation;
the assertion that the sum is a chain map and a unit is required only
after reduction modulo two.  End identifications, including any tensor
automorphisms, are fixed before composing the maps.
\end{theorem}

The proof occupies this section.  After fixing the complexes and index
conventions, we construct the
handle families and their determinant data.  We then establish the
perturbation and matching results needed to compute every boundary
face of the interval and pentagon counts.

\subsection{Complexes and index conventions}

\begin{lemma}[Reducibles at the ordinary ends]
\label{lem:ordinary:central-ends}
At a trivializable-determinant rational end occurring here, every
reducible flat connection is central after a square-root normalization
of the determinant.  Its adjoint cohomology has
$(h^0,h^1)=(3,0)$.  The admissible determinant on $Y_0$ has no reducible
projectively flat connection.  Small perturbations fixing the central
connections preserve these properties.
\end{lemma}
\begin{proof}
The first homology of the rational ends is either zero or $\Z/2$.
A reducible $\SU(2)$ representation is Abelian, so factors through that
first homology and takes values in $\{1,-1\}$.  Its adjoint local system
is therefore the trivial rank-three system.  Rational homology-sphere
cohomology gives $h^0=3$ and $h^1=0$.  The two possible central lifts
at an order-two end remain distinct fixed-determinant flat connections;
their adjoint corrections coincide.

A reducible projectively flat connection on $Y_0$ would split as two
line connections with equal real first Chern class.  Its determinant
would consequently pair evenly with the integral surface $A$, contrary
to the choice of determinant.  At a central rational flat the Hessian
on the Coulomb slice is invertible, since $h^1=0$.  The implicit function
theorem, with the residual compact stabilizer acting on that slice,
therefore preserves its isolation when the perturbation vanishes there.
If new reducible critical points appeared for perturbations tending
to zero, compactness in dimension three would give a reducible flat
limit.  The preceding isolation, or the admissibility obstruction at
$Y_0$, excludes such a sequence.
\end{proof}

Thus the irreducible-only Floer construction applies to these ends.
We use the determinant-one version throughout
\cite[Section~7.1, arXiv version~2]{KronheimerMrowka2010}; the extension to
three-manifolds with two-torsion first homology is explicitly recorded
in \cite[Section~2, Theorem~2.3]{GhoshSivekZentner2023}.  After a nondegenerate
perturbation there are finitely many generators.  The differential
counts oriented index-one trajectories modulo translation.  The
central-contact estimate proved below excludes central strata from
the compactifications defining the differential and its square.
At $Y_\infty$ choose the unperturbed flat problem, so that
\begin{equation}\label{eq:ordinary:sphere-complex}
 C(Y_\infty)=0.
\end{equation}
We may forget the finite Floer grading when discussing linear
isomorphisms; whenever a moduli space is counted, its actual integer
index, rather than only its residue class, is specified.

\begin{convention}[Relative sectors and fixed-limit index]
\label{conv:ordinary:index}
The cobordism $C_j$ is the oriented knot trace from $Y_\infty$ to $Y_j$.
Paths of cobordisms are written in chronological order; their induced
operators compose from right to left.  The symbol $C_j$ for a
cobordism is distinct from the complex $C(Y_j)$.

On a completed piece, $i$ denotes the unframed ASD index with the
boundary representatives fixed.  Gauge transformations may approach
independent parallel stabilizers at the ends.  Normal remainders have
small positive exponential weights.  We retain all relative bundle
sectors compatible with the fixed determinant transitions, and select
the required actual total index.  A matching across an isolated flat
adds $h^0$ to the sum of the two fixed-limit indices.

At $L=S^1\times S^2$ with trivialized determinant, harmonic flat
parameters are included before residual conjugation.  At a
noncentral flat $(h^0,h^1)=(1,1)$, and at a central flat
$(h^0,h^1)=(3,3)$.  Varying one fixed asymptote adds $h^1$; matching two
fixed-asymptote pieces adds $h^0+h^1$.

Choose the ordinary integral determinant-line orientation systems,
with true-end orientations, homology orientations and integral
determinant lifts fixed.  These systems are orientable for the
determinant-one problem.  In any fixed relative bundle sector the
space of reference connections with fixed product-flat collars is
affine, hence connected; a comparison made by determinant-line
excision therefore extends by orientation transport with a constant
relative sign.  Gluing these comparisons is associative.  Changing
an internal framing degree changes the relative charge sector, so
a cap prescribed to have charge $1/4$ uses that prescribed framing
degree.  These orientation conventions are those of the ordinary
ASD/Floer construction \cite{DonaldsonKronheimer1990}; see in particular
\cite[Section~5.4, Proposition~5.11, and Section~5.6]{Donaldson2002}
for determinant-line gluing and determinant-one conventions.
\end{convention}

\begin{lemma}[Charge and index bookkeeping]\label{lem:ordinary:index}
Put $E=-p_1(\ad E)/4$, interpreted by Chern--Weil on the completed
connection.  Increasing relative charge by one increases $i$ by eight.
If a connected piece has $b_1=0$ and only central rational ends, then
\begin{equation}\label{eq:ordinary:central-index}
 i=8E-3(1+b^+).
\end{equation}
For replacements with identical boundary flats and fixed transitions,
including all matching terms, the index difference is
\begin{equation}\label{eq:ordinary:excision}
 \Delta i=8\Delta E-\frac32\Delta(\chi+\sigma).
\end{equation}
The nonnegative lower charge bounds used below persist for sufficiently
small perturbations.  On a product neck with a bounded Floer functional
perturbation, the required error is bounded independently of the neck
length.
\end{lemma}
\begin{proof}
The real ASD symbol contributes $-3(\chi+\sigma)/2$, and its bundle
term is $-2p_1(\ad E)=8E$.  APS boundary terms for a central flat are
those of three trivial real coefficients.  Under
Convention~\ref{conv:ordinary:index}, the scalar fixed-limit complex
has index $-1-b^+$; tensoring by the parallel adjoint fiber gives
\eqref{eq:ordinary:central-index}.  This also gives $-3$ for a flat
central cylinder, as required by the addition of its three matching
directions.  Excision cancels the identical boundary terms and proves
\eqref{eq:ordinary:excision}.  These are the usual ASD index formulas
with the unframed correction retained; see \cite{DonaldsonKronheimer1990}
and \cite[Equation (10), p.~24]{DaemiLidmanMillerEismeier2025}.

For unperturbed ASD fields, the trace-free charge is a nonnegative
multiple of the curvature norm squared.  In each application its
possible values lie in a specified discrete translate of $\Z$.
If a sequence of perturbations tending to zero violated the stated
lower bound, compactness and charge quantization would give an ASD
limit, with any lost charges positive, violating that same bound.
On a long product one uses the gradient identity instead of a bound
proportional to length.  With the normalization of Chern--Simons
absorbed in a fixed positive constant, a gradient trajectory satisfies
\[
 \int_{a}^{b}\|\dot a(t)\|^2\,dt
 =\operatorname{CS}(a(a))+h(a(a))
  -\operatorname{CS}(a(b))-h(a(b)).
\]
The perturbation contribution is bounded by $2\|h\|_\infty$.
On concatenated product segments the intermediate functional values
telescope.  Consequently smallness can be fixed before increasing
the product length.  This is also the energy estimate used when
discarding a nonconstant collar level.
\end{proof}

A regular count over a $p$-dimensional metric family, without
insertions, has $i=-p$.  Its one-dimensional boundary relation has
$i=1-p$.  An irreducible nonconstant trajectory has index at least one
before quotienting by translation.  A \emph{particle} is an ideal
curvature-concentration point with a positive integral charge.  A particle
of charge $q$ lowers
the field index by $8q$ and has at most four location parameters.
These conventions will be used even when a piece is filled only
virtually for an index comparison.

\subsection{The polygon model and its metric families}

The polygon and associahedral surgery families follow the constructions
of Culler--Daemi--Xie~\cite[Section~4.3 and Theorem~4.56,
arXiv version~1]{CullerDaemiXie2020} and
Daemi--Lidman--Miller Eismeier~\cite[Section~5]{DaemiLidmanMillerEismeier2025}.
We supply the local plumbing, determinant and boundary calculations
needed for the present ends and integral units.

\begin{lemma}[Polygon traces and plumbing cuts]\label{lem:ordinary:polygon}
A path of adjacent slopes has a smooth toric local model over a polygon.
Choose clockwise primitive slope vectors $w_0,\ldots,w_m$ with
$\det(w_i,w_{i+1})=-1$.  Collapse the circle of slope $w_i$ over the
corresponding side and the whole torus at each corner, and attach the
remaining boundary to the knot exterior times an interval.  An arc
from the exterior boundary to side $i$ gives the primary cut
$Z_i=Y_{w_i}$.  An arc between sides $i,j$, $j-i\ge2$, gives a cut
$M_{ij}$ bounding the linear plumbing associated with the intervening
sides.  Its spheres have consecutive intersections $+1$ and squares
$-d_k$, where
\begin{equation}\label{eq:ordinary:fan-relation}
 w_{k-1}+w_{k+1}=d_kw_k.
\end{equation}
If the two endpoint slopes agree, the plumbing boundary is
$L=S^1\times S^2$.  Its primitive radical is the image of the boundary
$S^2$.  Replacing this cap by $D=S^1\times D^3$ joins the two equal
collapsed sides and leaves the exterior identity cylinder.
\end{lemma}
\begin{proof}
At a corner the two adjacent slope circles form an integral basis.
The quotient is therefore the usual smooth two-disk corner model,
rather than an orbifold.  Along an interior side the change from
the preceding basis to the following basis reverses the tangent
circle parameter and has normal degree $-d_k$, proving
\eqref{eq:ordinary:fan-relation} and the self-intersection formula.
The common corner contributes intersection $+1$ to neighboring
spheres.  Trivializations along the exterior side give precisely
the usual two-handle trace, since $\infty\to j$ is $C_j$,
$j\to\infty$ is its orientation-reversed reverse, and
$r\to r+1$ adds a $(-1)$-framed meridian to $C_r$.  Removing a
side of square $-1$ is the corresponding corner blow-down.

The inverse image of an arc from the exterior side is the filled
three-manifold at that slope.  The inverse image of an arc between
two collapsed sides is the boundary of their intervening chain,
which proves the cut descriptions.  When the endpoint circles
are equal, their collapse gives a sphere, and the transverse
circle remains free; the boundary is $S^2\times S^1$.
The sphere along that arc is primitive, and its image in the
plumbing is the primitive null vector of the intersection
matrix.  Filling by $S^1\times D^3$ replaces the arc by a
half-disk with one straight collapsed edge.  This joins the equal
sides without changing the exterior gluing and proves the last
assertion.
\end{proof}

\begin{lemma}[Compatible associahedral metric families]
\label{lem:ordinary:metric-families}
A two-step path has an interval family with endpoint cuts
$Z_1,M_{02}$.  A three-step path has a pentagon family with faces
$Z_1,Z_2,M_{02},M_{13},M_{03}$.  A four-step path has the corresponding
three-dimensional associahedral family.  Each face carries the product
families on its split pieces, and compatible simultaneous cuts are
exactly the disjoint noncrossing arc systems.

Two consecutive final handles may instead be retained as a fixed
segment.  For the path $2,\infty,0,1,2$, this gives a pentagon with
faces
\begin{equation}\label{eq:ordinary:compressed-cycle}
 Z_1,\quad Z_2,\quad M_{02},\quad M_{04},\quad M_{14}
\end{equation}
in cyclic order.  The $M_{04}$ edge has endpoints $M_{02}$ and
$M_{14}$; no cut internal to the fixed segment is introduced.
\end{lemma}
\begin{proof}
Use collars of the polygon arcs of
Lemma~\ref{lem:ordinary:polygon}.  At shared slope sides, separate
their endpoints slightly in noncrossing order; at the exterior
side use distinct ordered slices.  Thus every compatible arc
system has disjoint collars.  Face-product coordinates of an
associahedron specify independent lengths for precisely these
collars.  Choose the metrics on completed pieces inductively
over faces, with fixed product metrics on the cross-sections:
round metrics on lens spaces and the product metric on $L$.
For finite lengths glue and round these collars.  Once prescribed
on a neighborhood of the boundary of a parameter face, the
metrics extend over its compact interior because positive-definite
symmetric forms form a convex cone.

For two steps, the only alternatives are splitting at the
intermediate filling or isolating the intervening sphere.
For three steps, the noncrossing arc description gives the five
listed faces and their corner incidences.  Keeping the last two
handles fixed simply treats their union as the third block;
the same three-block construction applies.  In the original
indices its cap cuts are $M_{02},M_{14},M_{04}$, giving
\eqref{eq:ordinary:compressed-cycle}.  Omitting the interior
cut of a fixed block leaves that block's data unchanged and
does not create another boundary parameter.  The construction
inserts only neck lengths; it introduces no twist coordinate.
\end{proof}

\begin{lemma}[Topological and determinant data for the unit families]
\label{lem:ordinary:unit-topology}
The single steps $r\to r+1$ for $r=0,1$, and their composite
$0\to1\to2$, have $H_1=0$ and $b^+=0$.  A determinant on a single
step with the prescribed end restrictions is determined modulo
twice an integral class; changing it within that parity changes
the map only by tensor identifications at the ends.

The pentagons proving the single-step units may be taken on
\begin{equation}\label{eq:ordinary:two-triples}
 s\to\infty\to r\to s,
 \qquad r\to s\to\infty\to r,
 \qquad s=r+1.
\end{equation}
Their hard caps have chain $[-1,-1]$, primitive radical
$P_1+P_2$, and determinant evaluations $(1,-1)$ up to a common
sign.  The compressed pentagon for $g_+$ has chain
$[-2,-1,-2]$, radical $S_l+2T_1+R$, and determinant evaluations
\begin{equation}\label{eq:ordinary:hard-determinants}
 c_0(S_l,T_1,R)=(1,-1,1),
 \qquad c_1(S_l,T_1,R)=(-1,0,1).
\end{equation}
All these full pentagon cobordisms and their interval subpieces
have $b^+\le1$.
\end{lemma}
\begin{proof}
An attached meridian normally generates the incoming surgery
group, so the single steps have trivial first homology.  Precede
$r\to r+1$ by $C_r$ and blow down the meridional corner.  The
result is $C_{r+1}$ blown up once, with orthogonal classes
$G,T$, $G^2=r+1$, $T^2=-1$.  The incoming trace class is
$G-T$, and the image of the step's intersection lattice is its
orthogonal complement, by the relative homology exact sequence.
For $r=0$ it is generated by $G-T$ and is null; for $r=1$ it
is generated by $G-2T$, of square $-2$.  In particular neither
step has positive rank.  Preceding $0\to1\to2$ by $C_0$
gives the twice blown-up $C_2$, with incoming null class
$G-T_1-T_2$.  Its orthogonal complement has no positive
subspace.  This proves the assertions about $b^+$ and, again
by normal generation, $H_1$.

The rank-one integral cohomology of a step tests the indicated
generator, and the end restrictions fix its parity.  Since
$H_1=0$ there is no torsion ambiguity in this statement.  Two
determinants of that parity differ by twice a line class.
Tensoring the rank-two bundle by that line identifies their
adjoint equations and induces the corresponding fixed
identifications on the end complexes.  The rank and
invertibility assertions about the map are invariant under
these end automorphisms.

For the first triple in \eqref{eq:ordinary:two-triples}, use
orthogonal classes $F_l,G,T$ with squares $-s,s,-1$ and put
\[
 P_1=G-T-F_l,\qquad P_2=T,
 \qquad c(F_l,G,T)=(0,0,-1).
\]
Their intersections give the chain $[-1,-1]$ and
$c(P_1,P_2)=(1,-1)$.  For the second triple, blow down
$P_1=T$ to $-C_r+C_r$ with classes $F_l,G$ of squares
$-r,r$, and put
\[
 P_2=G-F_l-T,\qquad
 c(F_l,G,T)=(1,1,1).
\]
Again the chain determinant is $(1,-1)$.  The intermediate
$-C_s$ class is $F_l+T$, with even determinant evaluation;
thus the order-two end, when present, has the required
trivializable restriction.  The zero-surgery restriction is
odd in the relevant coordinates.  These presentations have
at most one positive direction, as do their subpieces.

For the compressed path the full trace is $-C_2$ followed by
the twice blown-up $C_2$.  In the orthogonal basis
$F_l,G,T_1,T_2$ the squares are $-2,2,-1,-1$.  Set
\begin{equation}\label{eq:ordinary:three-chain}
 \begin{split}
 A&=G-T_1-T_2,\\
 S_l&=G-T_1-T_2-F_l,\qquad R=T_2-T_1,\\
 c_0(F_l,G,T_1,T_2)&=(0,0,-1,0),\qquad
 c_e=c_0+e\PD(S_l).
 \end{split}
\end{equation}
Direct intersection gives squares $-2,-1,-2$ and adjacent
intersections $1$.  The determinant evaluations are precisely
\eqref{eq:ordinary:hard-determinants}, and both annihilate
the radical.  On $-C_2+C_0$ they restrict to the data of
Theorem~\ref{thm:ordinary:units}(ii).  At $Z_2$ the remaining
piece has an allowed determinant for $f_+$.  The displayed
ambient diagonal form also proves the positive-rank bound.
\end{proof}

\subsection{Compatible perturbations}

The handle families require perturbations that agree on shared faces
and vary independently on separated pieces.  Interior equation
perturbations provide this freedom; at true Floer ends they must be
gradients of time-independent invariant functions.  We construct
these two types of perturbation separately.

\begin{lemma}[Local perturbations and fixed-flat evaluation]
\label{lem:ordinary:perturbations}
For the finite ordinary families used in this section one can choose
arbitrarily small compatible perturbations with the following
properties.
\begin{enumerate}[label=\textup{(\roman*)}]
\item Irreducible equations, together with any specified finite
collection of auxiliary evaluations, are transverse, using variations
supported in the same component.  On a split face, variations on
different components are independent.
\item At a specified harmonic reducible, perturbations can vanish in
value while prescribing arbitrary finite-rank complex-linear maps
from the off-diagonal kernel to the off-diagonal cokernel.  Thus a
complex index-zero normal operator can be made invertible, and a
complex index-one normal operator can be made surjective.
\item Evaluation of an irreducible field at an $S^1\times S^2$ end
is transverse to a local flat slice.  Fixing a central flat imposes
three real conditions in that slice.
\end{enumerate}
The equation perturbations are independent of auxiliary cut marks.
They may be taken cylindrical at designated ordinary product ends,
with exponentially decreasing tails at separated auxiliary ends.
\end{lemma}
\begin{proof}
Choose a base frame and finitely many based loops and paths in the
component.  Their holonomies and transports are equivariant under
the compact frame group.  At an irreducible field some finite set
has precisely the central stabilizer.  To see the corresponding
infinitesimal assertion, suppose that every loop variation is a
simultaneous infinitesimal conjugation.  Extend that infinitesimal
conjugating element along paths from the base point.  The loop
condition makes the extension independent of the path, and the
resulting section is an infinitesimal gauge transformation.  Thus
finite samples detect any prescribed finite-dimensional horizontal
tangent space.

In a slice of the sample space choose equivariant coefficient
functions, and transport arbitrary compactly supported output forms
from small balls.  To test universal surjectivity at a smooth base
configuration, first use the auxiliary linearization from
$W^{1,q}_{\delta_S}$ to $L^q_{\delta_S}$, with $q>4$.  When the
higher charts use a H\"older weight $\delta_H$, choose
$0<\delta_S<\delta_H$ in the same root-free interval of the
normal end operator.  Field values on paths
and at points are continuous in this domain; coefficients arising
from moving paths or finite parameters are evaluated at the smooth
base.  Vary auxiliary cut values independently, thereby killing the
finite-dimensional cut component of any cokernel functional.  Its
equation component is then an $L^{q'}$ dual functional, where
$q'=q/(q-1)$.  Smooth compactly supported outputs detect every
nonzero such functional, even when required to avoid finitely many
labeled particle locations, since they are dense in $L^q_{\delta_S}$.
No smoothness of an adjoint cokernel field is asserted; derivatives
of graph terms can have distributional adjoints.  A finite selection
of outputs spans the finite-dimensional cokernel.  Primal elliptic
bootstrap transfers this surjectivity to the higher-order charts:
a low-order solution of the linearized equation with higher-order
right-hand side gains the required regularity.  The fixed output
forms are smooth, while transport and sample linearizations preserve
the available H\"older order with the prescribed summable bounds;
scalar cutoff functionals multiply smooth profiles.  Sobolev
embedding first supplies positive H\"older regularity, the local
elliptic inverse gains one order, and iteration gives each higher
order.  The end estimates transfer the solution to weight
$\delta_H$ because there are no intervening indicial roots;
no equal-weight H\"older-to-Sobolev inclusion is used.
A countable dense library and Sard--Smale then give simultaneous
transversality for the countably many relative sectors and breaking
types in question.

At a reducible, retain noncentral sample data to distinguish its
parallel torus.  The off-diagonal tangent and target spaces carry
the same complex torus character.  Slice-linear equivariant
functions with zero value on commuting samples realize arbitrary
complex-linear maps between finite kernel and cokernel subspaces.
Finite-rank perturbation theory then gives the asserted invertibility
or surjectivity.  The value on the entire Abelian path remains zero.
Further irreducible variations are supported away from that path in
configuration space.

To vary the asymptote, extend a prescribed infinitesimal harmonic
flat variation over the end.  Its equation error is compactly
supported after a cutoff and is cancelled by the same interior
universal variations.  This proves transversality of the evaluation.
Near a central $S^1\times S^2$ flat the slice consists of the three
parameters $h\vartheta$, $h\in\mathfrak{su}(2)$, with the full
parallel group acting on $h$.  On the irreducible exterior the
condition $h=0$ has codimension three.  Quotienting the flat space
first to a conjugacy interval would obscure these three equations
and is not our convention.

One may implement the samples by thickened loops, or by graph
samples multiplied by small-ball curvature cutoffs.  In the latter
description the cutoff is zero if the graph enters a ball with
curvature above a fixed small-energy threshold.  Consequently a
particle disables a graph it meets; graph terms away from particles
converge in regular gauges.  Countably many terms with summable
smooth operator norms suffice.  At product faces extend each term
from its component and turn off incompatible terms.  Product
cutoffs and a relative generic extension over parameter interiors
preserve the face data.  Moving paths can be handled in successively
higher Sobolev or little-H\"older orders, so that the finite
differentiability required by Sard--Smale is available.  No mark
variable enters the equation coefficients.  This construction is
the ordinary, zero-spinor specialization of the more detailed
compactness-compatible construction in
Section~\ref{sec:analysis}.
\end{proof}

For actual Floer ends the perturbations in the last lemma must be
gradients of time-independent invariant functions.  The next lemma
supplies this additional requirement, rather than treating arbitrary
interior curvature perturbations as Floer perturbations.

\begin{lemma}[Gradient transversality on ordinary trajectories]
\label{lem:ordinary:gradient-transversality}
Small time-independent invariant thickened-holonomy functions, fixing
the central critical points, can make the irreducible critical points
and all irreducible trajectories needed here regular.  Their
gradients are bounded, and their Hessians along a compact smooth
configuration arc are bounded operators on $L^2$.
\end{lemma}
\begin{proof}
Average invariant functions of simultaneous based holonomies over
a disk of base points and submersively thickened paths.  The first
variation is a bounded transported insertion along the paths;
Cauchy--Schwarz followed by integration over the thickening gives
an $L^2$ bound for the Hessian.  Use weighted summable smooth norms
for a countable collection of such functions.  Thin-loop values
and first variations on a compact smooth arc are limits of these
averages.  The detection argument of
Lemma~\ref{lem:ordinary:perturbations} therefore supplies arbitrary
finite-dimensional differentials on an irreducible sample slice,
which proves universal transversality for the critical points.

We verify the corresponding assertion for trajectories.  In temporal
gauge write $\dot a=-\grad(\operatorname{CS}+h)(a)$.  Two smooth
solutions agreeing on a time slice agree everywhere.  Indeed, on
a compact interval their difference satisfies
\[
 \dot\xi+D\xi=B(t)\xi,
\]
where $D$ is time-independent and symmetric and $B(t)$ is bounded
on $L^2$.  On a nonzero interval set
$m=\langle D\xi,\xi\rangle/\|\xi\|^2$.  Differentiation and
completion of the variance square give
\[
 m'\le \tfrac12\|B\|^2,
 \qquad
 (\log\|\xi\|^2)'=-2m+O(\|B\|).
\]
The first inequality bounds $m$ above towards the right, and the
second then prevents the norm from reaching zero at a finite
right endpoint.  Apply the same argument after reversing time and
replacing $D$ by $-D$ to exclude a finite left endpoint.  This proves
uniqueness in both directions.

If a spatial gauge transformation stabilizes one slice, uniqueness
implies that it stabilizes the entire trajectory.  Thus every slice
of an irreducible trajectory is irreducible.  Comparison with a
critical trajectory shows that the velocity never vanishes and
neither limiting critical point is attained.  Comparison with a
time translate shows that repetition modulo spatial gauge would
make the positive energy of a nonconstant segment repeat
indefinitely, contradicting finite energy.  Hence the trajectory
is an injectively immersed curve in the spatial quotient.

Let $\eta$ annihilate the universal linearization, including all
additional gradient perturbations supported away from the limiting
critical points.  Compactly supported temporal-connection and
spatial-field variations imply
\begin{equation}\label{eq:ordinary:trajectory-adjoint}
 d_a^*\eta=0,
 \qquad -\dot\eta+\mathcal H_a\eta=0,
 \qquad \mathcal H_a=*d_a+\operatorname{Hess}h(a).
\end{equation}
A small compact segment of the immersed trajectory can be isolated
from the rest by finitely many sample functions: injectivity,
compactness up to its limiting critical points, and
finite-dimensional tangent detection give such a slice chart.
In that chart choose a function whose restriction to the curve is
zero but whose differential in any chosen horizontal normal
direction is arbitrary.  Pairing its gradient with $\eta$ forces
$\eta(t)=c(t)\dot a(t)$.  Next choose arbitrary functions along the
curve, with compact support in the chart.  The universal pairing
becomes $\int c(t)\frac{d}{dt}f(a(t))\,dt=0$, so $c$ is constant.
Thickened averages approximate all these tests, hence give the
same conclusion.

Differentiating the trajectory equation gives
$\ddot a=-\mathcal H_a\dot a$.  Combining this with
\eqref{eq:ordinary:trajectory-adjoint} and $\eta=c\dot a$ shows that,
if $c\ne0$, both $\mathcal H_a\dot a$ and $\ddot a$ vanish.  The
$L^2$ norm of the velocity would then be constant; finite energy
forces it to be zero.  Thus $c=0$, proving universal surjectivity.
Local slices, elliptic regularity and simultaneous Sard--Smale
complete the construction.  This argument asserts no regularity
for arbitrary reducible trajectories; those are controlled by
their index and charge instead.
\end{proof}

The same choices regularize the metric families of
Lemma~\ref{lem:ordinary:metric-families}.  Relative cutoffs extend the
auxiliary data from their prescribed product faces, and relative
generic extension over each parameter interior leaves those face
restrictions unchanged.

\subsection{Ordinary gluing and harmonic necks}

At an irreducible nondegenerate match the end operator has no zero
modes.  We first recall this gluing argument, then treat the harmonic
flat modes at the auxiliary cap boundaries.

\begin{lemma}[Ordinary irreducible gluing]\label{lem:caps:irreducible-gluing}
Suppose completed ordinary ASD pieces have matching isolated,
nondegenerate irreducible end connections, fixed determinant-preserving
transition identifications, and compatible cut data.  At a regular
matched cut configuration, gluing gives the usual local product chart,
including the neck-length parameters when these are variable.  In
particular an isolated regular matched configuration gives one nearby
solution at each sufficiently long prescribed neck.  The chart exhausts
solutions converging to that configuration and carries the coherent
determinant-line gluing orientation.
\end{lemma}
\begin{proof}
At a nondegenerate irreducible critical point the gauge-fixed tangential
operator has $h^0=h^1=0$, hence a spectral gap.  The fields and their
linearized variations converge exponentially on the ends.  Pregluing
therefore has exponentially small equation and cut errors.  Split off
the finite-dimensional tangent space of the regular matched problem;
on its complement the half-problem operators have bounded right
inverses.  Their cutoff splice gives a parametrix on the joined neck.
There is no tangential zero mode to match.  Increasing the cutoff
transition length makes its derivative error small, and exponential
end decay controls the remaining splice error.  A Neumann series
therefore gives a right inverse on that complement, uniformly in
sufficiently long necks.  The gauge Laplacian has the same property.
For an isolated matched cut configuration, index zero makes these
operators invertible after the indicated gauge fixing and cuts.
The contraction theorem corrects the pregluing uniquely; with a
nonzero matched tangent space it gives the asserted smooth local chart.
Conversely, exponential end estimates put any converging solution in
this slice and in the contraction neighborhood, giving exhaustion.
These are the gapped irreducible gluing estimates of ordinary instanton
Floer theory \cite{DonaldsonKronheimer1990,Donaldson2002}.

The center $\{\pm1\}$ acts trivially on the adjoint data and extends
over either piece, so its two representatives do not give two matches.
The only permitted matching isomorphisms preserve the fixed determinant
transition.  Pairing adjacent end orientation lines in cobordism order
gives the asserted orientation.  In the cap applications of
Section~\ref{sec:caps} all cut degrees are even.  The asymptotic cut data have the decay required by the same
parametrix.
This argument is for irreducible matches; it does not use the
auxiliary reducible-cap hypothesis of
Lemma~\ref{lem:ordinary:regular-gluing}.
\end{proof}

The analytic input at an $S^1\times S^2$ neck differs from that at
a rational sphere.  In particular its central flat is not a smooth
endpoint of a one-dimensional slice.  We retain its three harmonic
coordinates until the gauge quotient is taken.

\begin{lemma}[Decay with harmonic flat modes]
\label{lem:ordinary:neck-decay}
Give $L=S^1\times S^2$ the product metric, and let $\alpha$ be a
flat connection with parallel stabilizer $H$.  On a sufficiently
small-energy neck, modulo gauge, the spatial field has the form
\begin{equation}\label{eq:ordinary:harmonic-slice}
 a(t)=\alpha+h(t)\vartheta+b(t),
 \quad h(t)\in H^0_\alpha(L;\ad),
 \quad b(t)\perp H^1_\alpha(L;\ad),
 \quad d_\alpha^*b(t)=0.
\end{equation}
Here $\vartheta$ is the circle one-form and the temporal component
has zero parallel part.  For $[-T,T]$ and any sufficiently small
fixed decay rate $\delta>0$, the normal part satisfies estimates
of the form
\begin{equation}\label{eq:ordinary:normal-decay}
 \|b(t)\|_{k}
 \le C_k\bigl(B_-e^{-\delta(t+T)}
                  +B_+e^{-\delta(T-t)}\bigr)
       +C_k\int_{-T}^T e^{-\delta|t-s|}\rho_k(s)\,ds,
\end{equation}
where $B_\pm$ measure the end-strip normal data and $\rho_k$ is
the decreasing perturbation or background error.  The derivative
of $h$ is bounded by the same normal size and error.  Its total
drift is therefore bounded independently of $T$, and on middle
strips the field is exponentially close to a single, possibly
varying, middle flat.

These assertions include a central $\alpha$, in whose fixed slice
$h$ has three real coordinates.  They give matching up to $H$,
exponential convergence on completed ends, and convergence of
holonomies uniformly over compact smooth families of bounded
paths on the relevant strips.  At isolated lens flats the same
assertions hold with no harmonic spatial term.
\end{lemma}
\begin{proof}
For a product metric the harmonic one-forms are exactly
$h\vartheta$, with $h$ parallel.  Since
$\vartheta\wedge\vartheta=0$, every $\alpha+h\vartheta$ in this
slice is an actual flat connection, including at a central
$\alpha$.  The residual parallel group acts by conjugation on
$h$.  Time-dependent parallel rotations remove the parallel part
of the temporal connection.

The slice equation determines the remaining temporal component
$\varphi$ by $d_\alpha^*d_a$ on the orthogonal complement of
parallel sections.  This operator remains invertible for small
$h,b$, and its right-hand side vanishes when $b=0$.  Its estimates
therefore have the form
$\|\varphi\|_{k+1}\le C_k\epsilon\|b\|_k+C_k\rho_k$ after
shrinking the slice.  Projecting the ASD equation to the
co-closed normal space gives
\[
 (\partial_t+K_\alpha)b=\mathcal R,
 \qquad K_\alpha=*d_\alpha,
 \qquad
 \|\mathcal R\|_k\le C_k\epsilon\|b\|_k+C_k\rho_k.
\]
There is no forcing from a pure harmonic field.  Projection to
harmonic forms gives the analogous bound for $|\dot h|$.

On the normal space $K_\alpha$ is self-adjoint with a spectral
gap.  Propagate its positive spectral part from the left and its
negative part from the right.  The resulting Green kernels are
bounded by $Ce^{-\lambda|t-s|}$ for a gap $\lambda>0$.
Convolution with these kernels and absorption of
$C\epsilon\|b\|$ in a weight $\delta<\lambda$ prove
\eqref{eq:ordinary:normal-decay}.  Interior elliptic estimates
give the higher derivative versions.  Integration of the
harmonic derivative gives a bound by
$C\epsilon(B_-+B_+)+C\int\rho$, with no factor $T$.
Integrating only between a middle strip and $t=0$ gives the
exponential closeness to $\alpha+h(0)\vartheta$ claimed in the
statement.

To justify using one slice along a long small-energy stretch,
start near a cluster flat on a middle strip.  Local compactness
places every boundary strip of a maximal slice interval close
to some flat.  The normal estimate and the length-independent
harmonic drift bound keep the solution inside a slightly larger
fixed slice until a strip leaves the small-energy regime.
Shrinking the energy threshold first excludes such an exit in
the stretches under discussion.  This also avoids changing a
three-dimensional central slice into a one-dimensional
noncentral slice during the argument.

Any remaining non-small stretch produces a nonconstant level
or a particle by ordinary compactness.  The collar charge on
$S^1\times S^2$ is integral because Chern--Simons is constant
on its connected flat space, so a nonflat level costs at least
one charge.  After the finitely many possible losses have been
extracted, the preceding estimates apply to all remaining
stretches.  Smooth convergence on strips implies the holonomy
assertion by the transport ODE and uniform bounds on the
chosen path family.  For a lens flat $h^1=0$, and the same
proof uses only the gapped normal operator.
\end{proof}

\begin{lemma}[Regular matching at an auxiliary reducible end]
\label{lem:ordinary:regular-gluing}
Suppose two completed ordinary pieces meet at one of the auxiliary
flats above.  Assume that one piece is a cap whose parallel stabilizer
is the whole matching group $H$; it has one matching end, and its
other end conditions, if present, are decaying.  Assume that the
connection on the other piece is irreducible; prescribed unglued
true-end asymptotes are also irreducible.  Frame the cap at the
matching end and leave the other piece unframed.  Suppose the
matched half-problem, including all flat parameters, metric parameters
and cuts, is isolated and transverse.  Then for all sufficiently long
necks each such match gives exactly one solution modulo
determinant-one gauge.  Every convergent solution sufficiently close
to this broken configuration arises in this way.  The regular
operators and gauge-fixing operators have inverses with at most
polynomial growth in neck length.

The assertion also holds for successive removal of pendant caps
in a tree of regular cuts, provided each removal satisfies these
hypotheses with compatible data.  Orientations are the determinant-line gluing
orientations; there is no factor from ordering a splitting, from
the central subgroup $\{1,-1\}$, or from additional projective
gauge twists.
\end{lemma}
\begin{proof}
Preglue the halves with their common flat.  By
Lemma~\ref{lem:ordinary:neck-decay}, the equation and cut errors are
exponentially small.  The gauge-augmented linearization on each half
has bounded extended kernel fields whose limits consist of a
parallel temporal mode and a harmonic spatial variation.  All normal
remainders decay exponentially.

The square operator on the finite glued manifold, with its true
ends completed, is
\[
 D_T:
 W^{k+1,q}_\delta(\Omega^1(\ad E))\oplus\R^{p+a}
 \longrightarrow
 W^{k,q}_\delta(\Omega^0(\ad E)\oplus\Omega^+(\ad E))
       \oplus\R^b.
\]
Here $q>4$, positive weights occur only at the fixed true ends,
$p$ counts metric parameters, $a$ counts cut marks, and $b$ is the
cut codimension.  The gauge row is $d_A^*$.  The irreducible core
eliminates a global infinitesimal stabilizer.  Index addition gives
\[
 \ind D_T=i_N^{\rm fix}+i_V^{\rm fix}+h^0+h^1+p+a-b=0,
\]
the last equality being isolated regularity of the matched problem.
Indeed varying both flat slices adds $2h^1$, matching subtracts
$h^1$, and framing the cap adds $h^0$.  At a central $L$ flat,
both $h^0$ and $h^1$ are three, so all three spatial matching
equations are retained.  This index-zero assertion concerns the
full finite operator.  The cap's decaying augmented operator can
still have the parallel scalar cokernel of its gauge row, even
when its physical framed problem is regular.  Extended temporal
modes and the following flux calculation account for that
cokernel; framing alone does not remove it.

We first identify the zero modes in the matched kernel.  On the
stabilizer cap integrate the linearized gauge-slice equation against
a globally parallel stabilizer section $\xi$.  If
$a=\tau\,dt+\beta+o(1)$, with $\tau\in H^0_\alpha$ and
$\beta\in H^1_\alpha$, integration on a truncation $N_R$ gives
\[
 0=\int_{N_R}\langle d_A^*a,\xi\rangle
   =-\int_{L_R}\langle a(\nu),\xi\rangle+o(1).
\]
There is one possible nondecaying boundary flux; all other end
terms tend to zero.  Global parallel sections restrict onto the
whole end stabilizer, so the limit forces
$\langle\tau,\xi_\infty\rangle=0$ for every $\xi_\infty$, and
hence $\tau=0$.  The spatial mode $\beta$ has no normal contraction.
The same argument kills all three temporal modes at a central
flat.  Although framing restricts the allowed cap gauges, the
parallel sections remain legitimate test sections in this
integration identity.  Temporal matching gives the same conclusion on the other
half.  The remaining spatial flat modes, field variations,
parameters and cut variations solve exactly the matched half
linearization.  Its kernel is zero by the isolated transversality
hypothesis.  A residual infinitesimal gauge in augmented gauge is
also zero: integration of its gauge Laplacian gives its covariant
derivative norm, with zero end terms.  At true irreducible ends use
decaying gauge parameters.

Here is a useful quantitative version of this argument.  Use a
uniform local norm, the supremum of the elliptic norms on unit
strips, together with the finite parameter norm and the fixed
positive outer-end weights.  If no polynomial lower bound held,
one could find unit-norm fields and parameters on necks of length
$2T$ with linearized residual $o(T^{-1})$ (and, if necessary,
smaller than any fixed inverse power for higher-order norms).
Their limits on the two halves solve the homogeneous half
equations.  Projection to the neck zero modes shows that their
derivatives are $o(T^{-1})$ plus the exponentially decreasing
normal tails.  More explicitly, writing $z_T$ for that projection,
between distances $R$ from the two neck ends one has
\[
 |z_T(T-R)-z_T(-T+R)|
 \le CT\|D_Tx_T\|_{\mathrm{uloc}}+Ce^{-\delta R}.
\]
First let $T$ and then $R$ tend to infinity.  Thus the two limiting
constants agree.  The
previous paragraph makes both half limits and all parameters
zero.  A remaining unit norm translating down the neck would
limit to a bounded homogeneous full-cylinder solution.  The
spectral gap leaves only a zero mode, but its value is already
zero by the same matching estimate.  This contradicts local
elliptic compactness.  In particular a mode of amplitude one
spread over a neck, whose derivative is of order $T^{-1}$,
cannot invalidate the argument: it is excluded by the stated
residual bound in the uniform local norm.

The full operator has index zero, so this lower bound implies
injectivity and hence surjectivity, giving an inverse with polynomial norm.
Equivalently one constructs its right inverse from the inverses
on the matched half spaces and the finite-dimensional neck zero
modes.  The gauge Laplacian has the same polynomial bound by
the long-cylinder Poincar\'e estimate and the irreducible core
on the other side.  Exponential pregluing errors dominate these
polynomial losses, so the contraction theorem produces one
solution in a shrinking slice neighborhood.

For exhaustion, a convergent sequence has, on its middle strips,
an exponentially close \emph{variable} flat, by
Lemma~\ref{lem:ordinary:neck-decay}.  Cut using that same flat on
both halves.  In weak positive weights on the normal remainders,
and the ordinary norm on the flat parameter, the half data
converge and have exponentially small equation and matching
errors.  The injective regular half linearization improves this
closeness modulo gauge to exponential closeness.  Hence the
sequence lies in the uniqueness neighborhood just constructed.
At a central flat the allowed stabilizer rotation acts on the
small harmonic variation as well as on the frame.

For the stated tree, remove the pendant stabilizer caps in the
specified order and use product gluing.  At each removal the flux
argument has just one nondecaying cap end; this is stronger than
merely assuming that the cut graph has no cycles.
The cap stabilizer absorbs the matching $H$-rotations.  The
central gauge subgroup is already part of the gauge quotient,
and the fixed determinant transitions do not permit summing
over nonliftable $\SO(3)$ twists.  A change of integral framing
degree changes the relative charge and hence the actual index;
it belongs to a different sector, not to an additional copy of
the same zero-dimensional match.  Finally the usual APS/Floer
determinant-line identifications, associative under these
gluings, give the stated integral orientations.
\end{proof}

\subsection{Cap shifts and primary-face exclusions}

\begin{lemma}[The elementary cap shifts]\label{lem:ordinary:cap-shifts}
The isolated exceptional-sphere faces of the single-step pentagons
contribute at least two indices relative to a flat-ball filling.
The $M_{14}$ face of the compressed pentagon has the same lower
bound.  The odd-determinant $(-2)$-disk-bundle cap at $M_{02}$
has a unique projective boundary flat class, of trace zero and
$(h^0,h^1)=(1,0)$, and its relative contribution including matching is
\begin{equation}\label{eq:ordinary:minus-two-shift}
 \delta=8(E-\tfrac18),\qquad E\in\tfrac18+\Z.
\end{equation}
At $E=1/8$ its only generic framed zero-dimensional match is the
harmonic reducible, with multiplicity one.  The two determinant
choices at this face have identical exterior data and cancel
modulo two.
\end{lemma}
\begin{proof}
For a split bundle write $d$ for the difference of the summand
classes.  Its trace-free charge is $-d^2/4$.  On an odd
exceptional sphere the least such square is $-1$, so
$E\in1/4+\Z$ and its nonnegative minimum is $1/4$.
The exceptional disk bundle and the flat ball have the same
$\chi+\sigma=1$.  Equation~\eqref{eq:ordinary:excision},
including the separating $S^3$ matching, therefore gives
the shift $8E\ge2$.

The $M_{14}$ chain has matrix and inverse
\[
 Q=\begin{pmatrix}-1&1\\1&-2\end{pmatrix},
 \qquad Q^{-1}=\begin{pmatrix}-2&-1\\-1&-1\end{pmatrix}.
\]
Its two determinant evaluation vectors are $(-1,1)$ and
$(0,1)$; both have square $-1$.  The cap is unimodular, its
link is $S^3$, and again $\chi+\sigma=1$.
Consequently $E\in1/4+\Z$ and the same shift is at least two.

For the $(-2)$ disk bundle an odd difference evaluates to one
on the section and has square $-1/2$, giving reference charge
$1/8$.  Its boundary is $\RP^3$.  The projective flat
representation with this determinant obstruction is represented,
after square-root normalization, by
$\diag(i,-i)$ on its generator; it is the unique such class.
Its stabilizer is a torus and $h^1=0$.

To count the charge-$1/8$ cap, we first show that its framed normal
index is zero.  Compactify the resolution to the ruled surface $\mathbb F_2$ and contract its
exceptional section to obtain $\mathbb P(1,1,2)$.  Let $T$
be a fiber class on $\mathbb F_2$.  It has degree one on the
exceptional section and corresponds to $\mathcal O(1)$ on
the common smooth complement.  The real off-diagonal ASD
symbol is complex; on a closed K\"ahler surface its complex
index is
$-\chi(\mathcal O(T))-\chi(\mathcal O(-T))$.
On $\mathbb F_2$ these fiber twists are pullbacks of
$\mathcal O(1),\mathcal O(-1)$ on the base, with no higher
direct image, giving Euler characteristics $2,0$.  On
$\mathbb P(1,1,2)$ the weight-one monomials give
$h^0(\mathcal O(1))=2$.  The monomial decomposition of the
three-coordinate \v{C}ech complex has no degree-one cohomology,
and its degree-two cohomology is dual to weighted polynomials
of degree $-k-4$ for $\mathcal O(k)$.  For $k=1,-1$ these
spaces vanish.  Thus there is no higher cohomology for
$\mathcal O(1)$ and no cohomology for $\mathcal O(-1)$.
Thus its two Euler characteristics are also $2,0$.
Excision on their common complex bundle complement has zero
normal index difference.  The cone ball has framed index
zero, as follows from the flat-ball operator restricted to
the invariant fields, so the framed normal index on the
resolution cap is zero.  Equivalently, the invariant flat-ball
complex has no harmonic one-forms or self-dual two-forms,
and framing removes its parallel gauge contribution.
The tangential framed operator has
index zero and no kernel or cokernel because $H^1=b^+=0$.
Unframing subtracts its one-dimensional stabilizer, and
matching adds that dimension back.  This gives zero total
shift at $1/8$.  All other relative sectors change charge
by an integer and index by eight, proving
\eqref{eq:ordinary:minus-two-shift}.  The complement
identifications here are complex bundle identifications;
matching just real Chern forms would not suffice for this
index excision.

At the reference charge the tangential harmonic solution is
unique.  Lemma~\ref{lem:ordinary:perturbations} makes the
index-zero complex normal operator invertible without changing
that solution.  A framed irreducible would carry a positive
dimensional orbit of the boundary torus, contradicting a
regular zero-dimensional framed count.  Additional levels or
particles would increase charge and are unavailable at the
minimum.  Finally adding $\PD(S_l)$ changes the determinant
evaluation by $-2$ on this cap, hence is tensoring by a line
of degree $-1$.  This tensor operation preserves the adjoint
equation, and its order-two action at the boundary preserves
the trace-zero class.  Choose tensor-related cap perturbations
and the same exterior perturbations.  The two minimum
contributions then agree and cancel over $\F_2$.
\end{proof}

\begin{lemma}[The primary central-group estimate]
\label{lem:ordinary:central-cost}
In a broken ordinary configuration, let a maximal consecutive group
consist of $q$ reducible pieces or levels with only central rational
contacts.  Suppose its pieces have $b_1=0$, and denote by
$b^+_{\rm group}$ the sum of their positive ranks.  Its index
contribution, including all internal matches and its two exterior
central matches, is at least
\begin{equation}\label{eq:ordinary:group-cost}
 3(1-b^+_{\rm group}).
\end{equation}
In particular this contribution is at least three when the group
has positive rank zero.  A direct central match between two
irreducible pieces adds three.  The estimate remains valid when
the isolated minimum odd $(-2)$ cap is virtually filled by its
reference cap.
\end{lemma}
\begin{proof}
Apply \eqref{eq:ordinary:central-index} to every group member
and use its nonnegative charge.  Their total is at least
$-3q-3b^+_{\rm group}$.  There are $q-1$ internal contacts
and two exterior contacts, each contributing three.  Adding
them gives \eqref{eq:ordinary:group-cost}.  Direct matching
uses the same stabilizer dimension.  For a group reaching a
prescribed irreducible true end, the exterior contact is
realized by a nonconstant trajectory from the central state
to that irreducible; this trajectory is not included in the
group and has index at least one.  No reducer can meet the
admissible end $Y_0$.  The last assertion follows from the
zero framed shift of the reference filling in
\eqref{eq:ordinary:minus-two-shift}; its removal changes
neither the exterior index convention nor the lower bound.
\end{proof}

\begin{lemma}[Compactifications in the unit dimensions]
\label{lem:ordinary:low-dimensional-compactness}
For the rigid maps, intervals and pentagons above, all moduli spaces
needed in dimensions zero and one have no particles and no unwanted
central contacts at primary cuts or true ends.  Primary $S^3$ faces
are empty.  Exceptional-sphere faces and the $M_{14}$ face are empty
in these dimensions.  The only possible contribution at the
compressed $M_{02}$ face is the paired minimum contribution of
Lemma~\ref{lem:ordinary:cap-shifts}, away from the corners.
All other primary boundary strata are ordinary chain breakings or
the stated compositions.
\end{lemma}
\begin{proof}
The positive subspaces of split main pieces inject orthogonally
into the unsplit form, so their total positive rank is at most
one.  If $p'$ parameters remain on irreducible pieces, their
total index is at least $-p'$ by regularity.  We now check each
possible number of primary cuts; this avoids supposing that
a non-admissible primary face behaves like an admissible one.

For a single rigid $b^+=0$ map every forbidden central group
has the surplus three in
Lemma~\ref{lem:ordinary:central-cost}.  The chain dimensions
are zero and one, so the surplus excludes it.  The same
argument permits an ordinary one-parameter continuation and
stretching the middle of $f_+$: even subtracting that one
parameter leaves a surplus two.

For an interval $p=1$ the relevant indices are $-1,0$; for
a pentagon $p=2$ they are $-2,-1$.  With no primary cut,
a central group of positive rank zero costs at least three,
leaving a lower total index $3-p$ and hence exceeding the
required indices.  If a central group uses the available
positive rank, it must contain the whole uncut cobordism
piece: a group confined to true-end cylinder levels has
positive rank zero.  It then requires two nonconstant
irreducible end trajectories.  They cost at least two in
total; there is no irreducible cobordism piece from which
to subtract the metric parameter count.

At one primary cut a group using the positive rank contains
a main piece incident to at least one true end and hence
requires at least one such trajectory.  At most $p-1$
parameters remain on the other pieces.  Its total is
therefore at least $1-(p-1)$, which is positive for the
interval and at least zero for the pentagon.  Both exceed
the needed boundary indices.  Other central groups again
have surplus three.  At two primary cuts of a pentagon
no metric parameters remain.  Every group contribution is
nonnegative and every irreducible piece has nonnegative
index, whereas the boundary problem has index $-1$.
This excludes its central configurations as well.  Since
$S^3$ has no irreducible generator, its primary faces are
empty in all these cases.

For an exceptional or $M_{14}$ cap, virtual filling leaves
the preceding nonnegative group bounds and adds at least
two indices.  At most one parameter remains on the other
pieces of a pentagon face, and none on an interval face.
Thus these faces cannot occur.  At $M_{02}$ the minimum
reference filling has shift zero.  It can occur at a
rigid exterior parameter, as in
Lemma~\ref{lem:ordinary:cap-shifts}; its corners impose
the additional parameter loss and are excluded by the
same primary and exceptional bounds.  A higher sector
there adds at least eight indices.

For completeness, particle positions do not defeat these
inequalities.  A lost charge $q\ge1$ lowers field index
by $8q$ and supplies at most $4q$ position parameters.
Thus the dimension of an ideal stratum is at most the
unbroken expected dimension minus $4q$, before any
additional face loss.  It is negative in dimensions at
most one.  If a graph perturbation or auxiliary evaluation
is disabled by the particle, retain its position as a
parameter and use the componentwise variations of
Lemma~\ref{lem:ordinary:perturbations}; no insertion of
positive degree occurs in these unit families to restore
the lost dimension.  The same argument applies to
particles on levels.  Charge discreteness and the bounded
gradient energy identity prevent accumulation of further
nonconstant levels.

Consequently regular compactness leaves precisely the
ordinary irreducible chain breaks and the stated
compositions.  Notice that reducible main components
joining two nondecaying primary ends have been
\emph{excluded} here by dimension.  We never apply
Lemma~\ref{lem:ordinary:regular-gluing} to such a component;
all its stabilizer absorbers in the following counts are
single-ended auxiliary caps.
\end{proof}

\subsection{The hard cap and its restriction degree}

At the remaining outer face, replacing the sphere chain by
$S^1\times D^3$ leaves an identity cylinder.  We will show that its
exterior contributes the continuation map and that the original cap
supplies an odd number of matches with each generic noncentral
boundary flat, summed over the two determinants in the compressed case.

\begin{lemma}[Flat replacement and the continuation exterior]
\label{lem:ordinary:flat-replacement}
At the outer cut $M_{03}$ of a single-step pentagon, or $M_{04}$
of the compressed pentagon, write $N$ for the sphere-chain cap
and $V$ for its exterior.  Replace $N$ by $D=S^1\times D^3$,
keeping the determinant transitions fixed.  The filled cobordism
is the appropriate identity cylinder, possibly with a fixed
change of end identification.  If $i_{\rm filled}$ denotes its
index, then
\begin{equation}\label{eq:ordinary:hard-index}
 i_{\rm total}=i_{\rm filled}+8E_N-3,
 \qquad i_{\rm filled}=i_V^{\rm fix}+h^1,
 \qquad E_N\in\tfrac14+\Z.
\end{equation}
For the pentagon boundary problem $i_{\rm total}=-1$, the only
possible exterior is regular and irreducible, has
$i_{\rm filled}=0$, and has a noncentral boundary flat.
Necessarily $E_N=1/4$.  Its count with the flat replacement cap
is the continuation isomorphism.
\end{lemma}
\begin{proof}
Both chain determinants annihilate their primitive radical, so
restrict trivially to $L$ and extend over $D$.  The last assertion
of Lemma~\ref{lem:ordinary:polygon} identifies the filled
cobordism.  Leave the replacement cap unperturbed.  A flat on
$L$ extends uniquely to $D$ with its chosen boundary
representative: both fundamental groups are generated by the
same circle.  Its fixed-limit unframed index is $-h^0$.
After framing it has a regular family of flats mapping
identically to the boundary flat slice.  Thus fixed-limit
addition gives
$i_V^{\rm fix}-h^0+(h^0+h^1)=i_V^{\rm fix}+h^1$.

For the two-sphere chain $\chi(N)=3$, $\sigma(N)=-1$;
for the three-sphere chain $\chi(N)=4$, $\sigma(N)=-2$.
In both cases $\chi(N)+\sigma(N)=2$, whereas
$\chi(D)+\sigma(D)=0$.  Excision therefore gives the first
identity of \eqref{eq:ordinary:hard-index}.  The split
reference differences in the two caps have square $-1$,
as computed explicitly in
Lemma~\ref{lem:ordinary:harmonic-caps} below.  Their charge
is $1/4$.  Chern--Simons is constant on the flat space of
$L$, and fixing the determinant and relative transitions
makes changes of the remaining bundle degree integral.
Hence $E_N\in1/4+\Z$, including charge on extracted hard-neck
levels.  Nonnegativity gives $E_N=1/4+k$ with $k\ge0$.
At total index $-1$, \eqref{eq:ordinary:hard-index} becomes
$i_{\rm filled}=-8k$.

The exterior has no remaining metric parameter; the one
parameter on this face belongs to the cap.  Its moduli space
with variable asymptote therefore has dimension
$i_V^{\rm fix}+h^1=i_{\rm filled}$.  The virtual filling
has positive rank zero, so
Lemma~\ref{lem:ordinary:central-cost} excludes true central
breaks there; an admissible end also excludes reducibles.
A central $L$-restriction imposes the full three conditions
of Lemma~\ref{lem:ordinary:perturbations}.  Thus generic data
leave only regular irreducibles of dimension zero, with a
noncentral restriction, and force $k=0$.

To identify their count, stretch the same $L$ neck in the
filled cylinder.  A nonflat $D$-side sector or collar level
costs at least one integral charge, lowering the remaining
regular index by eight.  Particles are excluded by the same
index loss.  The only limiting replacement fields are
therefore the regular flat family on $D$ just described.
Lemma~\ref{lem:ordinary:regular-gluing} applies: the global
parallel stabilizer on $D$ is the entire end stabilizer,
and $D$ has one matching end.  Its unique gluing identifies
this count with the usual continuation count.  Continuation
is a chain homotopy equivalence over $\Z$ with coherent
orientations.  A fixed tensor or collar identification can
compose it with an automorphism, which is still an
isomorphism.  In the compressed pentagon the exterior and
its data are common to the two determinant choices.
\end{proof}

\begin{lemma}[Minimum-charge harmonic caps]
\label{lem:ordinary:harmonic-caps}
Let $N$ be either hard cap of
Lemma~\ref{lem:ordinary:unit-topology}, with one of its specified
determinants $c$.  At charge $1/4$ the only possible unordered
split difference is $\{c,-c\}$.  For each metric of its cap
interval there is a unique Abelian ASD connection modulo gauge
with the ordered difference $d=c$, allowing the boundary flat
to vary.  These connections form a harmonic path.

At a fixed noncentral boundary flat the cap index is $-2$.
Its tangential unframed index is $-2$ and its complex normal
index is zero.  If a noncentral exterior angle is a transverse
value of the harmonic path, the framed cap with its one metric
parameter has one regular contribution at each crossing.
There is no second contribution from exchanging its summands
or from a residual boundary framing.
\end{lemma}
\begin{proof}
For the chain $[-1,-1]$ a class annihilating the radical has
evaluations $(a,-a)$ and square $-a^2$.  The determinant
parity makes $a$ odd, so charge $1/4$ forces $a=\pm1$.
For $[-2,-1,-2]$, write the difference evaluations as
$(d_l,d_m,d_r)$.  The radical condition and square are
\begin{equation}\label{eq:ordinary:three-chain-charge}
 d_l+2d_m+d_r=0,
 \qquad -\frac{d^2}{4}=\frac{d_l^2+d_r^2}{8}.
\end{equation}
The first and last evaluations are odd.  Charge $1/4$
forces $d_l,d_r\in\{1,-1\}$.  When $e=0$, $d_m$ is
odd, so $d_l=d_r$ and the possibilities are
$\pm(1,-1,1)$.  When $e=1$, $d_m$ is even, so
$d_l=-d_r$ and the possibilities are $\pm(-1,0,1)$.
This proves both the uniqueness assertion and $c^2=-1$.

The class $d$ restricts trivially in real cohomology to
$L$, hence lies in the image of compactly supported
cohomology.  Cylindrical-end Abelian Hodge theory gives
its $L^2$ harmonic representative.  The intersection form
on this image is negative definite and $b^+=0$, so that
representative is anti-self-dual.  Integrating it gives
an Abelian connection with the prescribed integral
summand class; its limiting flat is determined by this
connection.  The fixed determinant supplies the other
summand.  There is no missing self-dual obstruction when
the harmonic boundary parameter is allowed to vary.

For uniqueness, the difference of two such curvatures
is exact, anti-self-dual, and decays.  Stokes' theorem on
long truncations gives zero for its square; the flat
limiting difference contributes no boundary curvature.
An anti-self-dual exact difference of square zero
vanishes.  The plumbing has $H^1(N;\R)=0$ and is simply
connected, so a remaining flat difference of line
connections is gauge.  Smooth dependence of the inverse
Abelian operator on the metric produces the harmonic
path.  The same argument supplies the uniqueness of the
minimum odd $(-2)$ reference used in
Lemma~\ref{lem:ordinary:cap-shifts}.

At a noncentral flat the replacement cap has index $-1$.
Its difference from $N$ is $8(1/4)-3=-1$, so the fixed
cap index is $-2$.  In the diagonal scalar complex there
is one parallel stabilizer and one boundary-flux
obstruction when the flat is fixed; $H^1(N)=b^+(N)=0$
leaves no other tangential kernel or obstruction.
Thus the tangential index is $-2$.  Subtracting it from
the full real index leaves zero, so the complex normal
index is zero.  Framing restores one real dimension;
a transverse metric crossing of the prescribed flat
removes the remaining flux obstruction.  Derivative
variations from Lemma~\ref{lem:ordinary:perturbations}
make the complex normal operator invertible while
leaving the harmonic path fixed.  This gives a regular
zero-dimensional framed cap problem.  A putative
irreducible cap, unframed and at this fixed angle, would
have expected dimension $-2+1=-1$ and is absent for
generic data.

An unordered splitting is counted once.  Choose $d=c$
to order its two lines.  At a fixed noncentral boundary
representative its actual normalized eigenvalue, not
merely its conjugacy class, prescribes which boundary
eigenline is assigned to this ordered line.  Swapping
both lines describes the same unordered connection and
does not create a new boundary choice.  The parallel
torus of the split cap absorbs all matching rotations.
This is exactly the single-ended stabilizer situation
of Lemma~\ref{lem:ordinary:regular-gluing}, so there is
no additional multiplicity.
\end{proof}

\begin{lemma}[Normalized eigenholonomy and the odd restriction degree]
\label{lem:ordinary:holonomy-degree}
For the $[-1,-1]$ cap, the harmonic path has odd mod-two
intersection number with every generic noncentral conjugacy
class of boundary flats.  For the $[-2,-1,-2]$ cap the sum of
those intersection numbers over the two determinant choices
is odd.
\end{lemma}
\begin{proof}
Let $\beta$ generate the free circle of $L$, and choose
an integral spanning surface or chain $\Delta$ for $\beta$
in $N$.  Fix the determinant trivialization $\tau$ at
$L$, compatible throughout the family.  With curvature
normalized to integral periods, Abelian Stokes' theorem
for the ordered line of difference $d$ gives its
normalized eigenholonomy as
\begin{equation}\label{eq:ordinary:eigenholonomy}
 u(d,c)=(-1)^{\langle c,(\Delta,\tau)\rangle}
              \exp\left(\pi i\int_\Delta d_{\rm harm}\right),
\end{equation}
up to inversion.  Indeed the ordered summand curvature
is $(F_{\det}+d_{\rm harm})/2$.  Exponentiating its
integral over $\Delta$ gives the exponential in
\eqref{eq:ordinary:eigenholonomy}; the relative integral
of $F_{\det}$ gives the displayed half-determinant sign.
The sign cannot be discarded when determinants are
compared.  Inversion changes the ordering but not the
conjugacy class or the mod-two degree.

The primitive radical is the image of the boundary
sphere by Lemma~\ref{lem:ordinary:polygon}.  Boundary
duality says that a spanning chain of the free boundary
circle pairs once, up to sign, with that sphere.
Consequently
\begin{equation}\label{eq:ordinary:radical-pairing}
 \Delta\cdot(P_1+P_2)=1\pmod2
 \quad\hbox{or}\quad
 \Delta\cdot(S_l+2T_1+R)=1\pmod2,
\end{equation}
as appropriate.

We justify the curvature localizations used at the
ends of the harmonic paths.  Across a rational
plumbing link there are no real harmonic one-forms.
The linear version of the gapped estimate in
Lemma~\ref{lem:ordinary:neck-decay} gives exponential
decay for the harmonic curvature on the stretched
collar.  Hodge compactness on the fixed cores and
uniqueness of the harmonic representative identify its
limits with the prescribed piecewise representatives.
Their tail integrals tend to zero.  Extending compactly
supported representatives by zero therefore computes
the limiting integral over $\Delta$ by its relative
intersection pairing.  This is a linear Hodge
calculation; it does not require nonlinear gluing at
a higher-energy stratum.

For $[-1,-1]$, the two ends isolate the exceptional
spheres $P_1$ and $P_2$.  For the same ordered difference
$c=(1,-1)$ their limiting localized representatives are
respectively $-\PD(P_1)$ and $\PD(P_2)$.  Each yields
a central boundary value, because its integral over
$\Delta$ is integral.  Their quotient in
\eqref{eq:ordinary:eigenholonomy} is
$(-1)^{\Delta\cdot(P_1+P_2)}=-1$ by
\eqref{eq:ordinary:radical-pairing}.  The two endpoints
are therefore the opposite central values.

For $[-2,-1,-2]$, first consider the end isolating
$S_l$, with cap $B=\nu S_l$ and rational link
$Q=\partial B=\RP^3$.  Let $p$ be the integral relative
Thom class of $S_l$, represented by the line $\mathcal O(S_l)$
with its canonical nonzero section $s$ off $S_l$.
For the determinant lines choose
$D_1=D_0\otimes\mathcal O(S_l)$ and
$\tau_1=\tau_0\otimes s|_L$, as required by the fixed
outside identification.  Thus their relative first Chern
classes differ by $p$, not by an unspecified lift of its
absolute class.  In particular the two half-determinant
signs differ by $(-1)^k$, where
$k=\langle p,[\Delta]\rangle=\Delta\cdot S_l$.

The common remaining curvature can be seen explicitly.
In rational cohomology put $w=c_0+p/2$.  Then
\[
 p^2=-2,\qquad c_0\cdot p=1,\qquad
 w\cdot p=0,\qquad w^2=-\tfrac12,\qquad
 c_e=w+(e-\tfrac12)p.
\]
Under the $Q$ stretch, linear Hodge splitting localizes
$\pm p/2$ on $B$ and the same class $w$ on the complement.
Each part has charge $1/8$.  Only the \emph{difference}
of the two harmonic curvatures localizes entirely on $B$.
Its limit there is the harmonic representative of $p$,
and its limit on the complement is zero.  A compact
Thom representative and this harmonic representative
differ by an exact form whose primitive decays modulo
an exact boundary form, since $H^1(Q;\R)=0$.
Split $\Delta$ across $Q$ and apply Stokes' theorem;
the neck tails have vanishing integrals.  Consequently
the integral over $\Delta$ of the harmonic difference
tends to $k$.  The exponential factor in
\eqref{eq:ordinary:eigenholonomy} therefore changes by
$\exp(\pi i k)=(-1)^k$, cancelling the relative
determinant sign.
Thus the two paths have the same boundary value at
this end; that common value need not be central.

At the other end the unimodular cap on $T_1,R$ carries
the entire charge.  Its two localized differences are
\begin{equation}\label{eq:ordinary:unimodular-localization}
 d_0=\PD(T_1),\qquad d_1=-\PD(T_1)-\PD(R).
\end{equation}
The evaluations of these classes on $S_l,T_1,R$
are exactly \eqref{eq:ordinary:hard-determinants}.
Choose $\Delta$ to avoid this cap.  This is possible
because its separating link is $S^3$: van Kampen
applied to the simply connected plumbing shows that
the complement is simply connected, so $\beta$ bounds
there.  The curvature integrals in
\eqref{eq:ordinary:eigenholonomy} are now zero.
Equation~\eqref{eq:ordinary:radical-pairing} reduces
to $\Delta\cdot S_l=1\pmod2$.  Hence the determinant
signs of the two paths differ by $-1$.  Their values
at this end are the opposite central values.

Identify the conjugacy space of flat $\SU(2)$
connections on $L$ with the closed angle interval
$[0,\pi]$.
Figure~\ref{fig:ordinary:cap-endpoints} records the endpoint data.
\begin{figure}[htbp]
\centering
\begin{tikzpicture}[font=\small,
 sphere/.style={circle,draw,inner sep=1.7pt},
 path/.style={-{Stealth[length=1.6mm]},thick}]
 \node at (1,1.9) {$[-1,-1]$};
 \node[sphere,label=above:$-1$,label=below:$P_1$] (p1) at (0.3,1.1) {};
 \node[sphere,label=above:$-1$,label=below:$P_2$] (p2) at (1.7,1.1) {};
 \draw (p1)--(p2);
 \node (zero) at (4.2,1.1) {$0$};
 \node (pi) at (7.5,1.1) {$\pi$};
 \draw[path] (zero)--(pi);
 \node[align=center] at (5.85,0.4) {$P_1$ and $P_2$ isolations\\(either order)};
 \node at (1,-0.6) {$[-2,-1,-2]$};
 \node[sphere,label=above:$-2$,label=below:$S_l$] (sl) at (0,-1.4) {};
 \node[sphere,label=above:$-1$,label=below:$T_1$] (t1) at (1,-1.4) {};
 \node[sphere,label=above:$-2$,label=below:$R$] (r) at (2,-1.4) {};
 \draw (sl)--(t1)--(r);
 \node (common) at (4.2,-1.4) {$\theta_*$};
 \node (z0) at (7.5,-0.95) {$0$};
 \node (zpi) at (7.5,-1.85) {$\pi$};
 \draw[path] (common)--(z0);
 \draw[path] (common)--(zpi);
 \node[align=center] at (4.2,-2.55) {$S_l$ isolation};
 \node[align=center] at (7.5,-2.55) {$T_1,R$ isolation};
\end{tikzpicture}
\caption{Cap chains and their harmonic-path endpoint values in the
flat conjugacy interval.  The two compressed paths share a value
$\theta_*\in[0,\pi]$ and end at opposite central values.  The branches
are not assigned determinant labels.  Only the endpoint data are shown;
no monotonicity of the paths is asserted.}
\label{fig:ordinary:cap-endpoints}
\end{figure}

The preimage parity of a generic interior
angle under any path is determined by which side of
that angle contains each endpoint.  For the first
cap its two endpoints are $0,\pi$, so the parity is
one.  For the second cap the contributions of the
common endpoint cancel between the two paths, while
the other endpoints are $0,\pi$; their total parity
is again one.  Choose the exterior angles away from
the finitely many relevant limiting values and
transverse to these paths.  Such a choice is allowed
by Lemma~\ref{lem:ordinary:perturbations}.  Small
normal regularizations preserving the Abelian paths
do not change the parity.
\end{proof}

\begin{lemma}[No additional ends of the minimum cap matches]
\label{lem:ordinary:cap-end-compactness}
For a generic noncentral exterior angle, the cap matches counted
in Lemmas~\ref{lem:ordinary:harmonic-caps} and
\ref{lem:ordinary:holonomy-degree} form a compact zero-dimensional
set.  The two types of cap-interval endpoint give no extra
contribution.  Gluing them to the continuation exterior gives
exactly its count modulo two, summed over both determinants
in the compressed case.
\end{lemma}
\begin{proof}
At exceptional-sphere or unimodular-cap isolation, the isolated
part already carries charge $1/4$.  No charge remains on
its complement.  The limiting boundary angle is therefore
one of the localized harmonic endpoint values computed in
Lemma~\ref{lem:ordinary:holonomy-degree}, and a generic
exterior angle avoids it.

At the $S_l$ isolation in the compressed case, the
trace-zero $(-2)$ reference cap has its minimum charge.
After filling that piece, the other part has fixed
index $-2$ and no parameter.  It has no regular
irreducible solution.  A split solution extends
through the traceless link by the unique minimum
split filling and again has unordered difference
$\{c_e,-c_e\}$ by
Lemma~\ref{lem:ordinary:harmonic-caps}.  It therefore
has the already computed common endpoint angle,
also avoided by the exterior.  Additional charge
losses are impossible at total charge $1/4$;
particles have integral positive charge, and collar
levels between the same prescribed flat classes
also have positive integral charge.

All remaining matches occur at transverse interior
angles.  Their tangential and normal regularity and
multiplicity one were proved in
Lemma~\ref{lem:ordinary:harmonic-caps}.  The cap has
one matching end and absorbs its full parallel
stabilizer.  Lemma~\ref{lem:ordinary:regular-gluing}
therefore both constructs all nearby solutions and
exhausts every convergent sequence.  The odd degree
of Lemma~\ref{lem:ordinary:holonomy-degree} supplies
one copy modulo two of each continuation exterior.
No gluing at any excluded higher-energy cap or at
a reducible main connector has been invoked.
\end{proof}

\subsection{Boundary identities and the bridge}

\begin{proof}[Proof of Theorem~\ref{thm:ordinary:units}]
First form the rigid maps on the two increasing steps.  Their
chain identities and the identification of the rigid count
on $0\to1\to2$ with their composition follow from the
$b^+=0$ case of
Lemma~\ref{lem:ordinary:low-dimensional-compactness}, the irreducible
gluing of Lemma~\ref{lem:caps:irreducible-gluing}, and determinant-line
orientations.  All collars used in this comparison are
fixed.  Denote one single-step map by $u_r$.

For its inverse comparison use both triples in
\eqref{eq:ordinary:two-triples}.  The complementary
interval on $s\to\infty\to r$ has no non-chain
boundary: its $S^3$ endpoint is empty and its exceptional
cap endpoint has shift at least two.  Thus its index
$-1$ count is a chain map.  The five faces of the two
pentagons are as follows; the same labels refer to
arcs in their respective ordered paths.
\begin{center}
\begin{tabular}{@{}lll@{}}
\toprule
Face & $s,\infty,r,s$ & $r,s,\infty,r$\\
\midrule
$Z_1$ & empty $S^3$ split & composition across $Y_s$\\
$Z_2$ & composition across $Y_r$ & empty $S^3$ split\\
$M_{02}$ & exceptional cap, empty & exceptional cap, empty\\
$M_{13}$ & exceptional cap, empty & exceptional cap, empty\\
$M_{03}$ & continuation on $Y_s$ & continuation on $Y_r$\\
\bottomrule
\end{tabular}
\end{center}
Here the empty statements are precisely
Lemma~\ref{lem:ordinary:low-dimensional-compactness};
the last row is
Lemmas~\ref{lem:ordinary:flat-replacement}--\ref{lem:ordinary:cap-end-compactness}.
No other face exists by
Lemma~\ref{lem:ordinary:metric-families}.

The compactified one-dimensional pentagon moduli
spaces, of actual index $-1$, consequently give,
over $\F_2$, chain homotopies of the form
\begin{equation}\label{eq:ordinary:single-step-identities}
 u_r h_1\simeq J_s,
 \qquad h_2\widetilde u_r\simeq J_r,
\end{equation}
where $J_s,J_r$ are continuation isomorphisms.  The
map $\widetilde u_r$ has the determinant induced by
the second pentagon; it is tensor-equivalent to
$u_r$, by Lemma~\ref{lem:ordinary:unit-topology}.
The first identity makes $u_r$ surjective on homology.
The second makes $\widetilde u_r$ injective, and hence
imposes the opposite inequality on the dimensions of
the two finite-dimensional homology groups.  They
have equal dimension.  Therefore $u_r$ is an
isomorphism.  This argument does not require the
two complementary interval maps or their determinant
transitions to be identical.  Tensor equivalence
also proves the assertion for every allowed
single-step determinant.  It follows that
$f_+=u_1u_0$ is a unit.

Now use the paired interval on $2\to\infty\to0$.
Its $S^3$ endpoint is empty.  At its other endpoint
the two minimum trace-zero cap contributions agree
and cancel modulo two by
Lemma~\ref{lem:ordinary:cap-shifts}; every higher
sector is excluded.  The remaining ends are exactly
the differential terms.  Thus the paired index
$-1$ count $g_+$ is a chain map modulo two.

The compressed pentagon on $2,\infty,0,1,2$ has
exactly the faces in
\eqref{eq:ordinary:compressed-cycle}.  We record
each term separately:
\begin{center}
\begin{tabular}{@{}ll@{}}
\toprule
Face & Contribution to the paired count\\
\midrule
$Z_1$ & empty primary $S^3$ split\\
$Z_2$ & $f_+g_+$\\
$M_{02}$ & equal trace-zero terms, cancelling modulo two\\
$M_{14}$ & positive cap shift, empty\\
$M_{04}$ & continuation $J_2$ by combined odd degree\\
\bottomrule
\end{tabular}
\end{center}
At the $M_{04}$ endpoints the only possible additional
cap degenerations are $M_{02}$ and $M_{14}$, and
Lemma~\ref{lem:ordinary:cap-end-compactness} has
already excluded them at a generic exterior angle.
The compact one-dimensional paired pentagon count
therefore gives
\begin{equation}\label{eq:ordinary:two-step-identity}
 f_+g_+\simeq J_2\qquad\text{over }\F_2.
\end{equation}
Since $f_+$ and $J_2$ are isomorphisms, so is $g_+$.

Apply the same constructions to the mirror knot
and reverse the order of the ends, identifying the
oppositely oriented three-manifolds.  The ordinary
counts dualize.  Negating a determinant, or changing
the sign of the flip sphere, changes the two choices
by tensoring with lines, with the latter line
canonically trivial off the sphere.  Their prescribed
restrictions are the ones in part (iii), and the
mod-two unit argument is unchanged.  This proves
the claims for $f_-$ and $g_-$.
\end{proof}

\begin{corollary}[The ordinary bridge]\label{cor:ordinary:bridge}
Fix an orientation-preserving diffeomorphism
$\phi:Y_{-2}\longrightarrow Y_2$, with fixed bundle and collar
identifications.  The stacked cobordism and its rigid count
\begin{equation}\label{eq:ordinary:bridge}
 T=f_-\,\phi^{-1}\,f_+:Y_0\longrightarrow Y_0
\end{equation}
define an integral chain map inducing a unit over $\F_2$ and
an isomorphism over $\Q$.  The cobordism has
\begin{equation}\label{eq:ordinary:bridge-topology}
 H_1(T;\Z)=0,\qquad \chi(T)=4,\qquad \sigma(T)=-2.
\end{equation}
Its two increasing-handle halves are simply connected.
\end{corollary}
\begin{proof}
The maps $f_\pm$ have coherent integral orientations,
and the fixed diffeomorphism induces an integral
chain isomorphism.  Stretching their rational seams
introduces only $b^+=0$ primary pieces.  The surplus
three of Lemma~\ref{lem:ordinary:central-cost}
excludes central contacts in the chain dimensions,
so the rigid count equals the displayed composition
up to integral chain homotopy.  Each factor is a
unit modulo two by Theorem~\ref{thm:ordinary:units}.

To obtain the rational assertion, take the mapping
cone of this map of finite free integral complexes.
Its homology modulo two vanishes.  The universal
coefficient exact sequence implies that its
integral homology has no free summand, since any
such summand would survive tensoring with $\F_2$.
Thus the cone is acyclic over $\Q$, proving the
rational isomorphism.  No integral invertibility
of $T$ on torsion is required.

Each meridional handle kills a normally generating
meridian in the incoming surgery group.  Hence
both handle halves are simply connected and their
gluing has $H_1=0$.  There are four two-handles,
giving $\chi=4$.  Each half has one negative
direction and no positive direction: its remaining
intersection direction comes from the zero-surgery
boundary and is null, as in
Lemma~\ref{lem:ordinary:unit-topology}.  Signature
additivity across the rational seam gives
$\sigma=-1-1=-2$.
\end{proof}
